% ==============================================================================
% 1. CHINESE LANGUAGE SUPPORT (macOS Optimized)
% ==============================================================================
% fontset=none prevents ctex from auto-loading fonts that might conflict
\usepackage[UTF8, fontset=none, scheme=plain]{ctex}

% Use -RawFeature={script=CJK} to stop the font engine from looking for 
% the CJK script tag, which silences the warnings.
\setCJKmainfont{Songti SC}[RawFeature={-script=CJK}]
\setCJKsansfont{Heiti SC}[RawFeature={-script=CJK}]
\setCJKmonofont{Kaiti SC}[RawFeature={-script=CJK}]

% Define the Kai family
\setCJKfamilyfont{kai}{Kaiti SC}[RawFeature={-script=CJK}]

% Ensure \kaishu is defined and points to our family
\providecommand{\kaishu}{\CJKfamily{kai}}
\renewcommand{\kaishu}{\CJKfamily{kai}}

% --- Set Kaiti as the Default Chinese Font ---
\AtBeginDocument{\kaishu}

% ==============================================================================
% 2. BEAMER THEME & VISUAL STYLE
% ==============================================================================
\usetheme{default}
\usecolortheme{orchid}
\useinnertheme{rounded}
\usefonttheme{professionalfonts} 


% Global color settings for Title and Frame Titles
\setbeamercolor{title}{fg=lightmidnightblue}
\setbeamercolor{frametitle}{fg=lightmidnightblue}

\setbeamerfont*{frametitle}{size=\normalsize,series=\bfseries}
\setbeamertemplate{navigation symbols}{} 
\setbeamertemplate{footline}[frame number]
\setbeamerfont{footnote}{size=\small}

% ==============================================================================
% 3. ESSENTIAL MATH & SYMBOLS
% ==============================================================================
\usepackage{amsmath, amssymb, amsfonts, mathrsfs}
\usepackage{cancel, siunitx}
\usepackage{wasysym, pifont}

\newcommand{\cmark}{\ding{51}}%
\newcommand{\xmark}{\ding{55}}%
\newcommand{\omark}{\ding{109}}%

\newcommand{\cxmark}{%
  \begingroup
  \setlength{\fboxsep}{0pt}%
  \makebox[1.1em][c]{%
    \raisebox{0pt}{\ding{51}}% ✔
    \kern-0.65em
    \raisebox{0pt}{\ding{55}}% ✗
  }%
  \endgroup
}

% ==============================================================================
% 4. TABLES & ALIGNMENT
% ==============================================================================
\usepackage{colortbl}
\usepackage{tabularx, booktabs, array, multirow, makecell, diagbox, threeparttable}

% ==============================================================================
% 5. GRAPHICS, COLOR & MULTIMEDIA
% ==============================================================================
\usepackage{tcolorbox, graphicx, caption}
\usepackage[absolute, overlay]{textpos}
%\usepackage{multimedia, animate}
% ==============================================================================
% 6. TIKZ & DRAWING (Updated: Removed 'snakes', added 'decorations')
% ==============================================================================
\usepackage{tikz}
\usetikzlibrary{
    calc, 
    matrix, 
    patterns, 
    backgrounds, 
    arrows, 
    arrows.meta, 
    shapes, 
    circuits, 
    circuits.logic.US, 
    chains, 
    fit, 
    intersections, 
    positioning,
    decorations.pathmorphing, % Replacement for wavy/zigzag snakes
    decorations.pathreplacing, % Replacement for brace snakes
    decorations.shapes          % Replacement for shape-based snakes
}

% --- Custom TikZ Styles ---
\tikzset{
    shapenode/.style     = {draw, rectangle, fill=none, auto, node distance=0em, font=\normalsize}, 
    roundnode/.style     = {draw, rectangle, rounded corners=0.5em, inner sep = 0.2em, node distance=0em, minimum height=1.5em},
    rectanglenode/.style = {draw, rectangle, inner sep = 0em, node distance=0em, minimum height=2em},
    textnode/.style      = {draw=none, fill=none, rectangle, minimum size=0cm, inner sep = 0em, auto, node distance=0em, font=\normalsize},
    smalltextnode/.style = {draw=none, fill=none, rectangle, minimum size=0.5cm, auto, node distance=0em, font=\small},
    circlenode/.style    = {circle, draw, auto, inner sep = 0em, node distance=0em}, 
    ellipsenode/.style   = {draw, ellipse, inner sep = 0em, node distance=0em, minimum height=0em, minimum width=0em},
    dotnode/.style       = {draw, fill=black, circle, auto, minimum size=0.2em, inner sep = 0em, node distance=0em},  
    connect/.style       = {black, ->},
    implies/.style       = {double, double equal sign distance, -implies}, 
    every text node part/.style={align=center}
}
% ==============================================================================
% 7. ALGORITHMS
% ==============================================================================
\usepackage[ruled,lined,linesnumbered]{algorithm2e}
\IncMargin{1em}
\SetNlSty{text}{}{:}

% ==============================================================================
% 8. TEXT ANNOTATION & LAYOUT
% ==============================================================================
\usepackage{soul}
\usepackage[normalem]{ulem}
\usepackage{framed, boxedminipage, url}
\usepackage[english]{babel}
\usepackage{cancel} % for math

% ==============================================================================
% 9. CUSTOM COLORS
% ==============================================================================
\definecolor{lightmidnightblue}{rgb}{0.2, 0.2, 0.7}
\definecolor{mycolor}{rgb}{51,51,178}
\definecolor{darkgreen}{rgb}{0.0,0.5,0.0}
\definecolor{orchid}{rgb}{0.855, 0.439, 0.839}

\newcommand{\blue}[1]{{\color{blue} #1}}
\newcommand{\white}[1]{{\color{white} #1}}
\newcommand{\red}[1]{{\color{red} #1}}
\newcommand{\cyan}[1]{{\color{cyan} #1}}
\newcommand{\green}[1]{{\color{green} #1}}
\newcommand{\darkgreen}[1]{{\color{darkgreen} #1}}
\newcommand{\orange}[1]{{\color{orange} #1}}
\newcommand{\yellow}[1]{{\color{yellow} #1}}
\newcommand{\magenta}[1]{{\color{magenta} #1}}
\newcommand{\gray}[1]{{\color{gray} #1}}
\newcommand{\lightgray}[1]{{\color{gray!20} #1}}
\newcommand{\purple}[1]{{\color{purple} #1}}
\newcommand{\mdblue}[1]{{\color{lightmidnightblue} #1}}

% ==============================================================================
% 10. CUSTOM MACROS & COMMANDS
% ==============================================================================
\let\oldcite=\cite                                                                
\renewcommand{\cite}[1]{\blue{\oldcite{#1}}}
\renewcommand\bibname{References}

\newcommand{\cmark}{\ding{51}}
\newcommand{\xmark}{\ding{55}}
\newcommand{\dash}{\rule{0.05em}{0pt}\rule[0.3em]{0.6em}{0.06em}\rule{0.05em}{0pt}}
\newcommand{\mydash}{\rule{0.05em}{0pt}\rule[0.35em]{0.5em}{0.05em}}
\newcommand{\define}{\stackrel{\text{def}}{=}}
\newcommand{\codecomment}[1]{{\color{blue} //#1}}
\newcommand{\splitline}{\darkgreen{\rule[0.25\baselineskip]{\textwidth}{0.6pt}}}

\newcommand{\redul}{\bgroup\markoverwith{\textcolor{red}{\rule[-0.5ex]{2pt}{0.4pt}}}\ULon}
\newcommand{\blueul}{\bgroup\markoverwith{\textcolor{blue}{\rule[-0.5ex]{2pt}{0.4pt}}}\ULon}
\newcommand{\greenul}{\bgroup\markoverwith{\textcolor{green}{\rule[-0.5ex]{2pt}{0.4pt}}}\ULon}
\newcommand{\darkgreenul}{\bgroup\markoverwith{\textcolor{darkgreen}{\rule[-0.5ex]{2pt}{0.4pt}}}\ULon}
\newcommand{\blackul}{\bgroup\markoverwith{\textcolor{black}{\rule[-0.5ex]{2pt}{0.4pt}}}\ULon}
\newcommand{\lightgrayul}{\bgroup\markoverwith{\textcolor{gray!20}{\rule[-0.5ex]{2pt}{0.4pt}}}\ULon}

\newcommand{\iO}{i\mathcal{O}}
\newcommand{\diO}{di\mathcal{O}}
\newcommand{\sample}{\xleftarrow{\textup{\tiny R}}}
\newcommand{\AdvA}{\mathsf{Adv}_\mathcal{A}(\lambda)}
\newcommand{\expect}{\mathbb{E}}
\newcommand{\entropy}{\mathsf{H}}
\newcommand{\minentropy}{\mathsf{H}_\infty}
\newcommand{\avminentropy}{\tilde{\mathsf{H}}_\infty}
\newcommand{\Phisb}{\Phi_\textup{brs}^\textup{srs}}
\newcommand{\Oracle}{\mathcal{O}}
\newcommand{\Osign}{\mathcal{O}_\mathsf{sign}}
\newcommand{\Odec}{\mathcal{O}_\mathsf{dec}}

% ==============================================================================
% 11. BEAMER TEMPLATE PATCHES
% ==============================================================================
\makeatletter
\newcommand\soulcolor{%
  \let\set@color\beamerorig@set@color
  \let\reset@color\beamerorig@reset@color}
\makeatother

\makeatletter
\newcommand{\footlineextra}[1]{\gdef\insertfootlineextra{#1}}
\newbox\footlineextrabox
\defbeamertemplate*{footline extra}{default}{
  \begin{beamercolorbox}[ht=2.25ex,dp=1ex,leftskip=\Gm@lmargin]{footline extra}
    \insertfootlineextra
  \end{beamercolorbox}
}
\addtobeamertemplate{footline}{%
  \setbox\footlineextrabox=\vbox{\usebeamertemplate*{footline extra}}
  \vskip -\ht\footlineextrabox
  \vskip -\dp\footlineextrabox
  \box\footlineextrabox%
}{}
\let\beamer@original@frame=\frame
\def\frame{\gdef\insertfootlineextra{}\beamer@original@frame}
\footlineextra{}
\makeatother

% ==============================================================================
% 12. THEOREMS & CUSTOM BLOCKS
% ==============================================================================
\setbeamertemplate{theorems}[numbered] 
\theoremstyle{plain}
\newtheorem{proposition}[theorem]{Proposition}

\newenvironment{noteblock}[0]{%
\setbeamercolor{block body}{bg=blue!10,fg=black}\begin{block}{}}{\end{block}}
\newenvironment{remarkblock}[0]{%
\setbeamercolor{block body}{bg=darkgreen!10,fg=black}\begin{block}{}}{\end{block}}
\newenvironment{refblock}[0]{%
\setbeamercolor{block body}{bg=gray!20,fg=black}\begin{block}{}}{\end{block}}
\newenvironment{claim}{\begin{trivlist} \item \textbf{Claim.}}{\end{trivlist}}

\newenvironment<>{assumption}[1]{%
  \setbeamercolor{block title}{bg=orange,fg=white}
  \setbeamercolor{block body}{bg=orange!10,fg=black}
  \begin{block}#2{#1}}{\end{block}}
\newenvironment<>{construction}[1]{%
  \setbeamercolor{block title}{bg=orchid,fg=white}
  \setbeamercolor{block body}{bg=orchid!20,fg=black}
  \begin{block}#2{#1}}{\end{block}}
\newenvironment<>{exercise}[1]{%
  \setbeamercolor{block title}{bg=darkgreen,fg=white}
  \setbeamercolor{block body}{bg=darkgreen!10,fg=black}
  \begin{block}#2{#1}}{\end{block}}
\newenvironment<>{example}[1]{%
  \setbeamercolor{block title}{bg=darkgreen,fg=white}
  \setbeamercolor{block body}{bg=darkgreen!10,fg=black}
  \begin{block}#2{#1}}{\end{block}}
\newenvironment<>{conclusion}[1]{%
  \setbeamercolor{block title}{bg=red,fg=white}
  \setbeamercolor{block body}{bg=red!10,fg=black}
  \begin{block}#2{#1}}{\end{block}}

\newenvironment<>{greenbox}[1]{\begin{tcolorbox}[colback=green!5,colframe=green!40!black]#1}{\end{tcolorbox}}
\newenvironment<>{bluebox}[1]{\begin{tcolorbox}[boxsep=0em,colback=white,colframe=blue]#1}{\end{tcolorbox}}
\newenvironment<>{redbox}[1]{\begin{tcolorbox}[boxsep=0em, colback=white,colframe=red]#1}{\end{tcolorbox}}
\newenvironment<>{titledredbox}[1][2]{\begin{tcolorbox}[boxsep=0em, colback=white,colframe=red, title=#1]#2}{\end{tcolorbox}}
\newenvironment<>{orangebox}[1]{\begin{tcolorbox}[colback=white,colframe=orange]#1}{\end{tcolorbox}}
\newenvironment<>{question}[1]{\begin{center}\itshape{#1}}{\end{center}}
\newenvironment<>{theoremblock}[1]{\begin{noteblock}\mdblue{Theorem:} #1}{\end{noteblock}}
\newenvironment<>{lemmablock}[1]{\begin{noteblock}\mdblue{Lemma:} #1}{\end{noteblock}}

\newcommand{\Cloud}[2][180]% [angle], content
{\begin{tikzpicture}[overlay]
    \node[align=center, draw, shading=ball, text=white, cloud callout, cloud puffs=17, cloud puff arc=140, 
      callout pointer segments=3, anchor=pointer, callout relative pointer={(#1:2 cm )}, aspect=4,scale=0.5] at (0.2ex,0.5ex) {#2};
\end{tikzpicture}}